% !TeX root = ../main.tex
% ---------------------------------------------------------------------------
% Final Linings, Waterproofing and Durability
% Scope extension: not derived from the RSE3010 weekly material.
% Key sources: ITA WG on waterproofing; EN 1992 and AS 3600; ITA/AITES report
% on fire in tunnels; fib Bulletin on tunnel durability.
% ---------------------------------------------------------------------------
\chapter{Final Linings, Waterproofing and Durability}
\label{ch:final-linings}


\section{Lining philosophies and design}
\label{sec:lining-philosophies-and-design}

\subsection{Lining philosophies}
\label{sec:lining-philosophies}
\scopenote{Single-shell, double-shell and composite linings; what each assumes
about the long-term contribution of the primary support.}

\subsection{Cast in-situ final lining design}
\label{sec:cast-insitu}
\scopenote{Formwork travellers and pour sequence, unreinforced versus
reinforced linings, load cases governing the permanent condition, and
construction joint detailing.}

\subsection{Composite action between linings}
\label{sec:composite-linings}
\scopenote{Bonded and debonded interfaces, load sharing between primary and
secondary linings, and the long-term stiffness and creep assumptions that
determine it.}

\section{Drainage philosophy and waterproofing}
\label{sec:drainage-philosophy-and-waterproofing}

\subsection{Drained and tanked systems}
\label{sec:drained-tanked-systems}
\scopenote{The lining consequence of the hydrogeological decision in
Chapter~\ref{ch:hydrogeology}: drainage layers and invert drains versus a
fully tanked structure resisting full head.}

\subsection{Sheet membrane waterproofing}
\label{sec:sheet-membranes}
\scopenote{PVC and TPO membranes, geotextile protection, compartmentalisation
with waterbars, re-injectable hoses, and welding quality control.}

\subsection{Sprayed waterproofing membranes}
\label{sec:sprayed-membranes}
\scopenote{Bonded sprayed membranes enabling composite shell construction,
substrate and curing requirements, and the current limits of the evidence base.}

\section{Drainage systems and maintainability}
\label{sec:drainage-systems}
\scopenote{Invert and side drains, sedimentation and calcite clogging,
flushing access, and designing drainage that can actually be maintained over
the asset life.}

\section{Durability and serviceability}
\label{sec:durability-and-serviceability}

\subsection{Concrete durability in the tunnel environment}
\label{sec:concrete-durability}
\scopenote{Sulfate and acid attack, chloride ingress, carbonation,
alkali-silica reaction and leaching; exposure classification and binder
selection for aggressive groundwater.}

\subsection{Crack control and serviceability}
\label{sec:crack-control}
\scopenote{Crack width limits for watertightness and durability, early thermal
cracking in thick pours, and shrinkage and restraint effects.}

\section{Fire resistance and design life}
\label{sec:fire-resistance-and-design-life}

\subsection{Fire resistance and spalling protection}
\label{sec:fire-protection}
\scopenote{Design fire curves (RWS, HC, ISO 834), explosive spalling
mechanisms, polypropylene fibre dosage and sacrificial boards; structural
performance requirements during and after fire.}

\subsection{Design life and specification}
\label{sec:design-life}
\scopenote{100- to 120-year design lives, cover requirements, the role of
maintenance assumptions, and specification for durability rather than strength
alone.}

\section{Internal fit-out and secondary structures}
\label{sec:fitout}
\scopenote{Road and track slabs, walkways, cross-passage doors, ventilation
structures and services, and their load and access interfaces with the lining.}
